\documentclass[10pt,a4paper]{amsart}
\usepackage[margin=19mm]{geometry}
\usepackage{amsmath,amssymb,mathtools}
\usepackage[scr=rsfs,cal=euler]{mathalfa}
\usepackage{xfrac}
\usepackage[unicode,hidelinks,bookmarks=false]{hyperref}
\newcommand{\ie}{i.e.\ }
\newcommand{\cf}{cf.\ }
\newcommand{\resp}{respectively\ }
\newcommand{\Subsection}{\subsection}
\providecommand{\nobreakdash}{\nobreak\hbox{-}}
\setlength{\emergencystretch}{2em}
\multlinegap0pt
\usepackage{tikz}
\usepackage{amssymb}
\usepackage[normalem]{ulem}
\usepackage[shortlabels]{enumitem}
\usepackage{xcolor}
\usepackage{mathtools}
\newtheorem{theorem}{Theorem}
\newtheorem{lemma}{Lemma}
\newcommand{\G}{\mathsf{G}}	% specific graph symbol
\def\R{\mathbb R} % Real numbers.
\newcommand{\V}{\mathsf{V}}	% set of vertices
\def\e{\mathsf{e}} 	% edge
\def\v{\mathsf{v}}	% vertex
\title{\bf Doubly Nonlinear Diffusion Equations on Metric Graphs}
\author{J. M. Maz\'on, J. Toledo}
\date{Source: arXiv:2604.12589v1 (CC BY 4.0)}
\begin{document}
\maketitle

\noindent\emph{Source and licence.} \bf Doubly Nonlinear Diffusion Equations on Metric Graphs. Version arXiv:2604.12589v1 (CC BY 4.0), distributed by arXiv under the Creative Commons Attribution 4.0 International licence, http://creativecommons.org/licenses/by/4.0/, as recorded in the arXiv licence metadata for this exact version and shown on the abstract page https://arxiv.org/abs/2604.12589v1. Full licence notice: \texttt{licenses/arxiv-2604.12589-CC-BY-4.0.txt}.\par\smallskip

\noindent\textit{Editorial scope.} Counted unit: the article's lemma lemma3 (source0.tex lines 935--991). The article's complete original proof of it is delivered with it (source0.tex lines 993--1141, one \texttt{proof} environment of 149 lines). No mathematics is written, repaired or re-derived: the statement and the proof are the article's own lines, in the article's own notation, and no retained line is substituted or re-flowed.  Retained uncounted as the source-local context the proof needs: lines 783--785; lines 912--916.  Nothing was omitted from any retained span, so the package carries no omission record.  No retained line contains a citation, so the wrapper carries no bibliography and no bibliography entry.  No retained line contains a figure environment, an image-inclusion command or drawing code, so no image is delivered and the package declares no asset dependency.  Editorial formatting only, in addition: an AMS \texttt{amsart} preamble that restates the article's own declarations and macros used by the retained lines, namely the environments \texttt{theorem}, \texttt{lemma} on separate counters, together with the standard packages those lines need.  The retained environments print in this document as Theorem 1, Lemma 1 in order of appearance, whereas the article prints the same items with the numbering of its own full document.  The complete native source of the article as distributed in the arXiv e-print for this exact version is bundled unchanged as \texttt{sources/198412/source0.tex}.
\begin{equation}\label{test1}
\int_0^{\ell} v =\int_0^{\ell} f  +  b + a,
\end{equation}

 \begin{theorem}\label{teocontprin001} Let  $\G$ be a  connected and compact metric graph.
  Let $1 < p_\e < \infty$ be and   $\gamma_\e:\R\to \R$ be a continuous and increasing function with $\gamma_\e(\mathbb{R})=\mathbb{R}$ and $\gamma_\e(0)=0$ for all~$\e \in E(\G)$.
   For $v_i$ weak solutions  to $(P_{\omega^i}^{f_i})$, $i=1,2$, we have  
$$\int_{\G}(v_1-v_2)^+\le \int_{\G}(f_1-f_2)^+ + \sum_{\v\in\V(\G)}(\omega^1_\v-\omega^2_\v)^+\,.$$
\end{theorem}

  \begin{lemma}\label{lemma3} Assume $\gamma:\mathbb{R}\to \mathbb{R}$ is a continuous and increasing function with $\gamma(\mathbb{R})=\mathbb{R}$ and $\gamma(0)=0$. We have the  statements:
 \begin{itemize}
 \item[1.]  Let $v_\epsilon=\gamma(u_\epsilon)$ be weak solution to problem ($P_{p,a,b+\epsilon}^{\gamma,f}$),  $\epsilon \in \mathbb{R}$. Then:
\\ (i)  For $\epsilon>0$, we have have that
 $$\begin{array}{l}\displaystyle 
 u_0 \leq u_\epsilon,
 \\ \\
 \displaystyle u_0(\ell) <u_\epsilon(\ell);
\end{array} $$ 
  for $\epsilon<0$, we have that 
  $$\begin{array}{l}\displaystyle u_\epsilon\le u_0,
   \\ \\
 \displaystyle 
u_\epsilon(\ell)<u_0(\ell).
\end{array} $$ 
 \\ (ii) Moreover, $\epsilon\mapsto u_\epsilon(x)$ is continuous in $\mathbb{R}$ for all fixed $x\in[0,\ell]$, and 
\begin{equation}\label{ran001}
\begin{array}{l}
\displaystyle \lim_{\epsilon \to  +\infty} u_\epsilon(\ell) =  +\infty,
 \\ \\
\displaystyle \lim_{\epsilon \to  -\infty} u_\epsilon(\ell) =  -\infty.
\end{array} \end{equation}


  \item[2.] Let $v_\epsilon=\gamma(u_\epsilon)$ be weak solution to problem ($P_{p,a+\epsilon,b}^{\gamma,f}$), with $\epsilon \in \mathbb{R}$. Then:
  \\ (i)  For $\epsilon>0$, we have have that
 $$\begin{array}{l}\displaystyle 
 u_0 \leq u_\epsilon,
 \\ \\
 \displaystyle u_0(0) <u_\epsilon(0);
\end{array} $$ 
  for $\epsilon<0$, we have that 
  $$\begin{array}{l}\displaystyle u_\epsilon\le u_0,
   \\ \\
\displaystyle 
u_\epsilon(0)<u_0(0).
\end{array} $$ 
 \\ (ii) Moreover, $\epsilon\mapsto u_\epsilon(x)$ is continuous in $\mathbb{R}$ for all fixed $x\in[0,\ell]$, and 
\begin{equation}\label{ran001dos001}
\begin{array}{l}
\displaystyle \lim_{\epsilon \to  +\infty} u_\epsilon(0) =  +\infty,
 \\
\\
\displaystyle \lim_{\epsilon \to  -\infty} u_\epsilon(0) =  -\infty.
\end{array} \end{equation}

\item[3.] Let   $v_\epsilon=\gamma(u_\epsilon)$ be weak solution of $(P_{p,a-\epsilon,b+\epsilon}^{\gamma,f})$ for $\epsilon\in \mathbb{R}$.
\\
(i) For $\epsilon>0$ we have
$$u_\epsilon(\ell)> u_0(\ell)$$ and
$$u_\epsilon(0)< u_0(0).$$
 For $\epsilon<0$ we have
$$u_\epsilon(\ell)< u_0(\ell),$$ and
$$u_\epsilon(0)> u_0(0).$$
(ii) Moreover, $\epsilon\mapsto u_\epsilon(\ell)$ and $\epsilon\mapsto u_\epsilon(0)$ are continuous in $\mathbb{R}$. 
 \end{itemize}
 \end{lemma}

 \begin{proof}  {\bf Proof of Point 1}.    Take  $\epsilon>0$.  The monotonicity  $ u_0 \leq u_\epsilon$ is consequence of Theorem~\ref{teocontprin001} and the monotonicity of $\gamma$. 
 Let us now proof that $u_0(\ell) <u_\epsilon(\ell)$. 
By the mass balance property~\eqref{test1} we have
  \begin{equation}\label{epsilon1}
\int_0^{\ell} (v_\epsilon - v_0) = \epsilon.
  \end{equation}
We also have 
$$v_0 = \left(|u_0'|^{p-2}u_0' \right)' + f,$$
and
  $$v_\epsilon = \left(|u_\epsilon'|^{p-2}u_\epsilon' \right)' + f. $$
Hence, for $0 \leq x \leq \ell$, having in mind the monotonicity and \eqref{epsilon1}, we have
$$0\le \int_0^x \left(  |u_\epsilon'|^{p-2}u_\epsilon'  -  |u_0'|^{p-2}u_0'   \right)'
= \int_0^x (v_\epsilon - v_0) \nearrow \epsilon \quad \hbox{as } x\nearrow \ell.$$
Hence,
$$0\le  (|u_\epsilon'|^{p-2}u_\epsilon') (x) -  (|u_0'|^{p-2}u_0') (x) - \left((|u_\epsilon'|^{p-2}u_\epsilon') (0) -  (|u_0'|^{p-2}u_0') (0) \right) \nearrow \epsilon \quad \hbox{as } x\nearrow \ell.$$
Now, 
$$(|u_\epsilon'|^{p-2}u_\epsilon') (0) = a \quad \hbox{and} \quad (|u_0'|^{p-2}u_0') (0) = a,$$
therefore, we have
\begin{equation}\label{17531510}
0 \leq (|u_\epsilon'|^{p-2}u_\epsilon') (x) -  (|u_0'|^{p-2}u_0') (x) \nearrow \epsilon \quad \hbox{as } x\nearrow \ell.
\end{equation}
Then, since the function $\varrho(r):= \vert r \vert^{p-2}r$ is   increasing for $p >1$, we obtain, from the first inequality in~\eqref{17531510},
$$u_0'(x) \leq u_\epsilon'(x),$$
 and integrating here, for $0\le x\le y\le \ell,$ 
$$u_0(y) - u_0(x)\leq u_\epsilon(y) -  u_\epsilon(x),$$
so,    
\begin{equation}\label{cry001}
0 \leq  u_\epsilon(x)  - u_0(x)\leq u_\epsilon(y) - u_0(y)\le u_\epsilon(\ell) - u_0(\ell).
\end{equation}
 Hence, if $u_\epsilon(\ell) - u_0(\ell)=0$, then $u_\epsilon=u_0$, which is false. Then we have $u_0(\ell) <u_\epsilon(\ell)$.

 The proof for $\epsilon<0$ is similar, and we have proved (i). 
 
 Let us now prove  (ii). By~\eqref{epsilon1}, there exists a  sequence $\epsilon_n\to 0$ such that $v_{\epsilon_n}\to v_0$ almost everywhere. Then, by~\eqref{cry001}, 
\begin{equation}\label{3.12.2}u_\epsilon(x)\to u_0(x)\quad\hbox{as }\epsilon\to 0^+,
\end{equation}
for any $x\in[0,\ell)$. 
Let us see that $$u_\epsilon(\ell)\to u_0(\ell)\quad\hbox{as }\epsilon\to 0.$$ 
From~\eqref{17531510} we also have: 
\begin{equation}\label{conv0003}
u_\epsilon'(x)\le \varrho^{-1}(\varrho(u_0'(x))+\epsilon).
\end{equation}
 Hence, integrating between $0$ and $\ell$,
 $$ u_\epsilon(\ell) - u_\epsilon(0)\le \int_0^\ell \varrho^{-1}(\varrho(u_0'(x))+\epsilon)dx,$$
 and therefore,
 \begin{equation}\label{12061610}
0\le u_\epsilon(\ell) - u_0(\ell)\le \int_0^\ell \varrho^{-1}(\varrho(u_0'(x))+\epsilon)dx+u_\epsilon(0)-u_0(\ell).
\end{equation}
By the  Monotone Convergence Theorem, 
 $$\lim_{\epsilon\to 0}\int_0^\ell \varrho^{-1}(\varrho(u_0'(x))+\epsilon)dx=\int_0^\ell u_0'(x)dx=u_0(\ell)-u_0(0),$$
we moreover have shown in~\eqref{3.12.2} that  $$u_\epsilon(0)\to u_0(0),$$   hence
 $$ 
 \displaystyle\lim_{\epsilon\to 0}\int_0^\ell \varrho^{-1}(\varrho(u_0'(x))+\epsilon)dx+u_\epsilon(0)-u_0(\ell)=0,
 $$
 then, from~\eqref{12061610},
 $$u_\epsilon(\ell)\to u_0(\ell)\quad\hbox{as }\epsilon\to 0^+.$$
 
The limit when  $\epsilon\to 0^-$ is similar.   Hence, we have  the continuity of $\epsilon\mapsto u_\epsilon(\ell)$ at $0$. Therefore,  fixed $x\in[0,\ell]$, having in mind~\eqref{cry001},   we have  the continuity of $\epsilon\mapsto u_\epsilon(x)$ at $0$. Now, it is easy to see that we have the continuity of $\epsilon\mapsto u_\epsilon(x)$ at any point~$\epsilon$.  
   

Let us finally prove the  statements in~\eqref{ran001}.   Take $\epsilon>0$.      By \eqref{epsilon1},
$$\int_0^\ell v_\epsilon(x)dx=\int_0^\ell v_0(x)dx+\epsilon,$$ and from here, there exists $0\le x_\epsilon\le\ell$ such that 
$$\lim_{\epsilon\to+\infty}v_\epsilon(x_\epsilon)=+\infty,$$
hence, 
$$\lim_{\epsilon\to+\infty}u_\epsilon(x_\epsilon)=+\infty.$$
Consequently,   since  by~\eqref{cry001},
$$  u_\epsilon(x_\epsilon) - u_0(x_\epsilon)\le u_\epsilon(\ell) - u_0(\ell),$$
and $u_0(x_\epsilon)$ is bounded, we get 
$$\lim_{\epsilon\to+\infty}u_\epsilon(\ell)=+\infty.$$
The other limit  follows in the same way.

The proof of {\bf Point 2.} is similar.

 Let us now prove {\bf Point 3.}  Take $\epsilon>0$. 
In this case   \begin{equation}\label{1746411new0}
\int_0^\ell (v_\epsilon-v_0)=0.
\end{equation}

 Since
$$(|u_\epsilon'|^{p-2}u_\epsilon')(\ell)-(|u_0'|^{p-2}u_0')(\ell)=\epsilon,$$
 and $x\mapsto (|u_\epsilon'|^{p-2}u_\epsilon')(x)-(|u_0'|^{p-2}u_0')(x)$ is continuous, we have that
$$(|u_\epsilon'|^{p-2}u_\epsilon')(x)-(|u_0'|^{p-2}u_0')(x)>0$$
on an interval $(x_\epsilon,\ell)$. Let us consider that this interval is the  the largest one (with $\ell$ fixed) for such property.  Let us now  consider the following two cases, (a) and (b).

\noindent (a) If $x_\epsilon=0$ then, since
  $$(|u_\epsilon'|^{p-2}u_\epsilon')(0)-(|u_0'|^{p-2}u_0')(0)=\epsilon,$$  
  we have that
$$(|u_\epsilon'|^{p-2}u_\epsilon')(x)-(|u_0'|^{p-2}u_0')(x)>0\quad\forall x\in[0,\ell].$$
Hence 
\begin{equation}\label{arre002new} u_0'(x) <u_\epsilon'(x)\quad\forall x\in[0,\ell],
\end{equation}
and integrating  
\begin{equation}\label{arre001new}   u_\epsilon(x)-u_0(x)< u_\epsilon(y)-u_0(y)\quad\forall 0\le x<y\le\ell.
\end{equation}
Now, if  $u_\epsilon(\ell)\le u_0(\ell)$,  by~\eqref{arre001new},
$$u_\epsilon(x)-u_0(x)\le 0\quad\forall 0\le x \le\ell,$$
and also, since $\gamma$ is increasing, 
$$v_\epsilon(x)-v_0(x)\le 0 \quad\forall 0\le x \le\ell.$$
Then, by~\eqref{1746411new0},
we have that $v_\epsilon=v_0$, which is impossible. Therefore 
$$u_\epsilon(\ell) >u_0(\ell).$$
And similarly we have that 
$$u_\epsilon(0) <u_0(0).$$
Let $\widehat x_\epsilon>0$  be such that $u_\epsilon(x)-u_0(x)>0$ for  $\widehat x_\epsilon<x\le \ell$
and  
\begin{equation}\label{mue02001}u_\epsilon(\widehat x_\epsilon)=u_0(\widehat x_\epsilon)
\end{equation}
Then, 
$$(|u_\epsilon'|^{p-2}u_\epsilon')(x)-(|u_0'|^{p-2}u_0')(x)-\left((|u_\epsilon'|^{p-2}u_\epsilon')(\widehat x_\epsilon)-(|u_0'|^{p-2}u_0')(\widehat x_\epsilon)\right)$$
$$= \int_{\widehat x_\epsilon}^x(v_\epsilon-v_0) \nearrow \epsilon-\left((|u_\epsilon'|^{p-2}u_\epsilon')(\widehat x_\epsilon)-(|u_0'|^{p-2}u_0')(\widehat x_\epsilon)\right) \quad \hbox{as } x\nearrow \ell.$$
Therefore,
$$(|u_\epsilon'|^{p-2}u_\epsilon')(x)-(|u_0'|^{p-2}u_0')(x)\le  \epsilon \quad \hbox{for } \widehat x_\epsilon<x<\ell.$$
From here, and using~\eqref{mue02001},
\begin{equation}\label{simcas01001}
0\le u_\epsilon(\ell)-u_0(\ell)\le \int_{\widehat x_\epsilon}^\ell\varrho^{-1}\left(\varrho(u_0'(x))+\epsilon\right)+u_0(\widehat x_\epsilon)-u_0(\ell).
\end{equation}


\noindent (b) If 
  $0<x_\epsilon<\ell$, then we have that
$$(|u_\epsilon'|^{p-2}u_\epsilon')(x_\epsilon)=(|u_0'|^{p-2}u_0')(x_\epsilon).$$
So, if
we apply the first point proved on this interval $(x_\epsilon,\ell)$,  we get that $u_0(x)\le u_\epsilon(x)$ in such interval and $u_0(\ell)< u_\epsilon(\ell)$; and we get that $u_0(x)\ge u_\epsilon(x)$ on the interval $(0,x_\epsilon)$ and $u_0(0)> u_\epsilon(0)$.
Observe that at point $x_\epsilon$, by continuity,
\begin{equation}\label{tkae001} u_\epsilon(x_\epsilon)=u_0(x_\epsilon).
\end{equation}
 Moreover, the proof of  the first point gives us the following fact  (using~\eqref{tkae001}): 
\begin{equation}\label{simcas01002}0\le u_\epsilon(\ell)-u_0(\ell)\le \int_{x_\epsilon}^\ell\varrho^{-1}\left(\varrho(u_0'(x))+\epsilon\right)+u_0(x_\epsilon)-u_0(\ell).
\end{equation}


 Observe that in both cases, (a) or (b), we obtain the similar inequalities~\eqref{simcas01001} and~\eqref{simcas01002}. Let us rewrite $x_\epsilon$ as $\widehat x_\epsilon$ in~\eqref{simcas01002}.

Then, like in the proof of the first point,  we get that
$$\lim_{\epsilon\to 0^+}u_\epsilon(\ell)=u_0(\ell).$$
 Indeed, we have that there exists $\widehat x_{\epsilon_n}\to x_0\in[0,\ell]$ for a sequence $\epsilon_n\to 0^+$. 
 Now,
$$
\int_{\widehat x_{\epsilon_n}}^\ell\varrho^{-1}\left(\varrho(u_0'(x))+\epsilon_n\right)=\int_{\widehat x_{\epsilon_n}}^{x_0}\varrho^{-1}\left(\varrho(u_0'(x))+\epsilon_n\right)+\int_{x_0}^\ell\varrho^{-1}\left(\varrho(u_0'(x))+\epsilon_n\right).
$$
Here, the fist integral on the right hand-side converges to $0$ as $\epsilon_n$ goes to $0$ since  the function $\varrho^{-1}\left(\varrho(u_0'(x))+\epsilon_n\right)$ is bounded uniformly in $n$, and, using this boundedness again,  the Dominate Convergence Theorem gives that the second integral on the right hand-side converges to $u_0(\ell)-u_0(x_0)$.  Then, we have
$$u_{\epsilon_n}(\ell)\to u_0(\ell)\quad\hbox{as }\epsilon_n\to 0,$$ 
and consequently, by the monotonicity,
$$\lim_{\epsilon\to 0^+}u_\epsilon(\ell)=u_0(\ell).$$


Similarly we get $\lim_{\epsilon\to 0^-}u_\epsilon(\ell)=u_0(\ell).$ 
Hence we have the  continuity of $\epsilon\mapsto u_\epsilon(\ell)$  at~$0$, and then at any point.
\end{proof}

\end{document}
